\documentclass[11pt]{article}
\usepackage[letterpaper,margin=1in]{geometry}
\usepackage{amsmath,amssymb,amsthm,mathtools,booktabs,tabularx,array,microtype,xurl}
\usepackage[T1]{fontenc}
\usepackage[colorlinks=true,linkcolor=blue,citecolor=blue,urlcolor=blue]{hyperref}
\newtheorem{theorem}{Theorem}[section]
\newtheorem{certtheorem}[theorem]{Computer-assisted theorem}
\newtheorem{proposition}[theorem]{Proposition}
\newtheorem{lemma}[theorem]{Lemma}
\newtheorem{corollary}[theorem]{Corollary}
\theoremstyle{remark}
\newtheorem{remark}[theorem]{Remark}
\newcommand{\R}{\mathbb R}
\newcommand{\C}{\mathbb C}
\newcommand{\eps}{\varepsilon}
\newcommand{\tstar}{t_*}
\newcommand{\Nzero}{N_0}
\newcommand{\Nstar}{N_*}
\newcommand{\Lsharp}{L_N^{\mathrm{sh}}}
\newcommand{\pos}[1]{\left(#1\right)_+}
\DeclareMathOperator{\Rea}{Re}
\DeclareMathOperator{\Log}{Log}
\DeclareMathOperator{\Capbd}{Cap}
\hypersetup{pdftitle={Sharp finite-sum bounds for the Riemann heat flow},
 pdfauthor={Drafted for Edward Baker},
 pdfsubject={Unreviewed working draft, 17 October 2026}}
\numberwithin{equation}{section}
\title{Sharp finite-sum bounds for the Riemann heat flow\\
\large Certified strip transfer, a finite window, and an infinite tail}
\author{Drafted for Edward Baker}
\date{17 October 2026\\\small Unreviewed working draft}
\begin{document}
\maketitle
\begin{abstract}
We study a normalized sharp lower bound for the paired finite sums in
Polymath's effective approximation to the Riemann heat flow. At heat time
$157/1000$, we certify monotonic transfer in the imaginary coordinate
throughout $[\sqrt{87677/2500000},\sqrt{343/500}]$ for four Euler-mollified
index ranges covering $690988\le N\le3840000$. The proof combines
active-branch right-derivative inequalities, exact analytic cases and
outward interval arithmetic over the actual integer parameters. A generic
signed-sum example shows why the old plain-weight transfer cannot simply
be inherited. One certified finite floor then gives a complete strip
margin exceeding $0.0017752269101384707$ in its natural cutoff window.
A separate eleven-term mollifier, with enlarged convolution head, gives
an infinite-tail margin exceeding $0.02303064$ from $N=3840000$ onward.
The remaining finite floors are not supplied, so the proposed bound
$\Lambda\le0.1745354$ does not follow from these results.
\end{abstract}

\noindent\textbf{Draft status.}
This draft assembles the results and execution records of round-one
Notes 01--02 \cite{note01,note02}. Published approximation results and the
historical initial-verification and barrier inputs remain explicit
imports. The computational checks are internal, LLM-assisted checks by the
same model; round-level cross-model review and specialist human review
remain pending. The authoring configuration and assistance are recorded
in the acknowledgements.

\section{Purpose and normalization}\label{sec:setup}

The Riemann heat flow in the normalization of Polymath \cite{polymath} is
\begin{equation}\label{eq:H}
 H_t(z)=\int_0^\infty e^{tu^2}\Phi(u)\cos(zu)\,du,
\end{equation}
where
\[
 \Phi(u)=\sum_{n\ge1}
 (2\pi^2n^4e^{9u}-3\pi n^2e^{5u})e^{-\pi n^2e^{4u}}.
\]
The present calculation concerns a quantitative finite-sum estimate at a
fixed positive time. Its main result is a strip-transfer theorem for a
sharp coefficient mass, not a criterion equivalent to the Riemann
hypothesis. The underlying sharp triangle inequality is inherited from
the Polymath method and the preceding audit \cite[Section 8.5]{polymath}\cite{audit3,sharpreview};
we give the native normalization and its proof because the application
requires the paired coefficients and their changing absolute-value branches.

Fix throughout
\begin{equation}\label{eq:parameters}
 \tstar=\frac{157}{1000},\quad
 y_0=\sqrt{\frac{87677}{2500000}},\quad
 y_1=\sqrt{\frac{343}{500}},\quad
 \Nzero=690988,\quad \Nstar=3840000.
\end{equation}
In particular $y_1^2=1-2\tstar$. The finite Euler prime sets are
specified by Table~\ref{tab:legs}. All statements over a leg quantify over
\emph{every integer} in its indicated range.

\begin{table}[htbp]
\centering
\begin{tabular}{lll}
\toprule
Leg & Integer indices $N$ & Euler prime set $P$\\
\midrule
0 & $690988\le N\le728999$ & $\{2,3,5,7,11\}$\\
1 & $729000\le N\le818999$ & $\{2,3,5,7\}$\\
2 & $819000\le N\le1027999$ & $\{2,3,5\}$\\
3 & $1028000\le N\le3840000$ & $\{2,3\}$\\
\bottomrule
\end{tabular}
\caption{The finite legs. The last natural index also lies in the tail.}
\label{tab:legs}
\end{table}

\subsection{The imported effective approximation}

For clarity we specify the functions used in the approximation. On
$\C\setminus(-\infty,1]$ use the principal logarithm and put
\begin{align}
 \mathcal M_0(s)&=\frac{s(s-1)\sqrt{2\pi}}{16}\pi^{-s/2}
 \exp\!\left[\left(\frac s2-\frac12\right)\Log\frac s2-\frac s2\right],\label{eq:M0}\\
 \alpha(s)&=\frac1{2s}+\frac1{s-1}+\frac12\Log\frac{s}{2\pi},
 &\mathcal M_t(s)&=e^{t\alpha(s)^2/4}\mathcal M_0(s).
 \label{eq:alpha}
\end{align}
For $z=x+iy$, set
\[
 s_+=\frac{1+y-ix}{2},\qquad s_-=\frac{1-y+ix}{2},\qquad
 B_t(z)=\mathcal M_t(s_+),
\]
\begin{align}
 s_*&=s_++\frac t2\alpha(s_+),
 &\gamma&=\frac{\mathcal M_t(s_-)}{\mathcal M_t(s_+)},\label{eq:sstar}\\
 \kappa&=\frac t2\bigl[\alpha(s_-)-\alpha(\overline{s_+})\bigr],
 &b_t(m)&=e^{(t/4)\log^2m}.\label{eq:kappa}
\end{align}
The natural cutoff and finite approximation are
\begin{align}
 N&=\left\lfloor\sqrt{\frac{x}{4\pi}+\frac t{16}}\right\rfloor,\label{eq:cutoff}\\
 f_t(x+iy)&=\sum_{m=1}^N b_t(m)m^{-s_*}
 +\gamma\sum_{m=1}^N m^y b_t(m)m^{-\overline{s_*}-\kappa}.
 \label{eq:f}
\end{align}
Theorem 1.3 of \cite{polymath}, for $0<t\le1/2$, $0\le y\le1$ and
$x\ge200$, supplies
\begin{equation}\label{eq:approx}
 \left|H_t/B_t-f_t\right|\le e_A+e_B+e_{C,0}.
\end{equation}
Its bounds for $|\gamma|$, $\Rea s_*$ and $|\kappa|$ are the inputs
used below. The carrier $B_t$ is nonzero: $s_+$ is off the real cut,
all factors in \eqref{eq:M0} are nonzero there, and the heat factor is an
exponential. We retain the full $e_{C,0}$ allowance. The error supplier uses
the conservative constant $10.50=3.58+6.92$ and denominator $x-12$,
as derived from Proposition 6.6(vi) in \cite{polymath} and audited in
\cite{audit4}. No residual from a different corrected center is inserted.

\subsection{Frozen windows and coefficients}

At $t=\tstar$ define
\begin{equation}\label{eq:window}
 q_N=N^2-\frac{\tstar}{16},\qquad x_N=4\pi q_N,
 \qquad W_N=[x_N,x_{N+1}).
\end{equation}
For $D=\prod_{p\in P}p$ and $d\mid D$, let
\begin{equation}\label{eq:coeff}
 \lambda_d=(-1)^{\omega(d)}\prod_{p\mid d}b_{\tstar}(p),\qquad
 \begin{aligned}
 B_{N,n}&=\sum_{\substack{d\mid D,\ d\mid n\\ n/d\le N}}
                  \lambda_db_{\tstar}(n/d),\\
 A_{N,n}(y)&=\sum_{\substack{d\mid D,\ d\mid n\\ n/d\le N}}
                  \lambda_db_{\tstar}(n/d)(n/d)^y .
 \end{aligned}
\end{equation}
Both leading coefficients are one. The frozen scalar bounds are
\begin{align}
 g_N(y)&=\exp\!\left(\frac y{50}-\frac y2\log q_N\right),\label{eq:g}\\
 \sigma_N(y)&=\frac{1+y}{2}+\frac{\tstar}{4}\log q_N
 -\frac{\tstar}{2x_N^2}
 \pos{1-3y+\frac{4y(1+y)}{x_N^2}},\label{eq:sigma}\\
 k_N(y)&=\frac{\tstar y}{2(x_N-6)},\qquad
 M_N(y)=\sum_{d\mid D}|\lambda_d|d^{-\sigma_N(y)},\label{eq:M}\\
 C_N(y)&=\sum_{m=2}^N b_{\tstar}(m)(m^{k_N(y)}-1)m^{y-\sigma_N(y)},
 \qquad K_N(y)=g_N(y)C_N(y).\label{eq:K}
\end{align}
The imported bounds imply throughout $W_N$
\begin{equation}\label{eq:freeze}
 |\gamma|\le g_N<1,\qquad \Rea s_*\ge\sigma_N,
 \qquad |\kappa|\le k_N.
\end{equation}
Indeed the first and third majorants decrease with $x$. In the lower
bound for $\Rea s_*$, the logarithm increases, and the product of
$x^{-2}$ with the nonnegative decreasing positive part decreases. This
also handles its kink. These directions require no monotonicity of the
final lower bound in independent scalar arguments.

\section{The sharp native binding}\label{sec:binding}

For real $B,A$ and $0\le g<1$, define
\begin{align}
 U_g(B,A)&=\max\!\left(\frac{|B-gA|}{1-g},\frac{|B+gA|}{1+g}\right)
 =\frac{|B-g^2A|+g|A-B|}{1-g^2},\label{eq:U}\\
 W_g(B,A)&=(1-g)U_g(B,A).
 \label{eq:weight}
\end{align}
The normalized coefficient mass and lower bound are
\begin{align}
 V_N(y)&=\sum_{n=2}^{DN}U_{g_N(y)}(B_{N,n},A_{N,n}(y))n^{-\sigma_N(y)},
 \label{eq:V}\\
 \Lsharp(y)&=\frac{(1-g_N(y))(1-V_N(y))}{M_N(y)}-K_N(y).
 \label{eq:Lsharp}
\end{align}

\begin{lemma}[Paired real coefficients]\label{lem:paired}
Let $r=|\Gamma|<1$. The real-linear map $T_\Gamma(w)=w+\Gamma\bar w$
satisfies $|T_\Gamma(w)|\ge(1-r)|w|$, and
\[
 T_\Gamma^{-1}(Bz+\Gamma A\bar z)
 =\frac{(B-r^2A)z+\Gamma(A-B)\bar z}{1-r^2}.
\]
In particular the modulus on the right is at most $U_r(B,A)|z|$.
For fixed real $A,B$, $U_r(B,A)$ is nondecreasing in $r$.
\end{lemma}
\begin{proof}
The inverse formula follows by substitution. The triangle inequality
and \eqref{eq:U} give the modulus estimate. For monotonicity, simultaneous
sign reversal reduces to $B\ge0$. For $0<r<1$ the three formulas for $U_r$
are
\[
 \frac{B-rA}{1-r},\quad \frac{B+rA}{1+r},\quad\frac{rA-B}{1-r}
\]
on $A\le B$, $B\le A\le B/r^2$, and $A\ge B/r^2$, respectively.
Their derivatives are $(B-A)/(1-r)^2$, $(A-B)/(1+r)^2$ and
$(A-B)/(1-r)^2$, all nonnegative on the corresponding branches.
The formulas agree at their boundaries and extend continuously to $r=0$.
\end{proof}

\begin{proposition}[Binding in the approximation's native units]\label{prop:binding}
For an index and prime set in Table~\ref{tab:legs} and $y\in[y_0,y_1]$,
if $\Lsharp(y)>0$, then
\[
 |f_{\tstar}(x+iy)|\ge\Lsharp(y)\qquad(x\in W_N).
\]
There is one division by $M_N$; the correction $K_N$ is not divided again.
\end{proposition}
\begin{proof}
Write $f=F_B+\gamma F_{A,\kappa}$ for the two sums in \eqref{eq:f},
and remove $m^{-\kappa}$ to define $F_{A,0}$. For
$E(s_*)=\prod_{p\in P}(1-b_{\tstar}(p)p^{-s_*})$, exact convolution gives
\[
 EF_B=\sum_{n=1}^{DN}B_{N,n}n^{-s_*},\qquad
 \overline E F_{A,0}=\sum_{n=1}^{DN}A_{N,n}n^{-\overline{s_*}}.
\]
Here $E\ne0$, since $\sigma_N>1.6>(\tstar/4)\log p$ makes each
Euler-factor perturbation have modulus less than one. Put
$\Gamma=\gamma E/\overline E$, so $|\Gamma|=|\gamma|=r$.
Lemma~\ref{lem:paired}, applied termwise with leading term $T_\Gamma(1)$,
yields
\[
 |E(F_B+\gamma F_{A,0})|
 \ge(1-r)\left(1-\sum_{n=2}^{DN}U_r(B_{N,n},A_{N,n})n^{-\Rea s_*}\right).
\]
If $V_N<1$, \eqref{eq:freeze} and monotonicity of $U_r$ imply that
this is at least $(1-g_N)(1-V_N)$. Also
\[
 |F_{A,\kappa}-F_{A,0}|
 \le\sum_{m=2}^N b_{\tstar}(m)(m^{k_N}-1)m^{y-\sigma_N}=C_N.
\]
Use $|E|\le M_N$ and subtract $g_NC_N$ after division. Finally,
$\Lsharp>0$ implies $V_N<1$, since $M_N>0$ and $K_N\ge0$.
\end{proof}

\begin{remark}
The weight $W_g=\max(|B-gA|,(1-g)|B+gA|/(1+g))$ need not increase with
$g$. The valid frozen-$g$ argument uses $U_g$ and positivity of $1-V_N$.
The improvement retains the paired real coefficients; a separate absolute
bound on the two sums would recover the plain triangle loss.
\end{remark}

\section{Right derivatives and the transfer obligation}\label{sec:dini}

For real $z,w$ set
\[
 \mathcal D(z,w)=
 \begin{cases}\operatorname{sgn}(z)w,&z\ne0,\\|w|,&z=0.\end{cases}
\]
It is the right derivative of the absolute value of a differentiable real
function with value $z$ and derivative $w$. We write $D^+$ for the upper
right Dini derivative. Fix $N,n$, abbreviate $L=\log n$, and set
\begin{equation}\label{eq:ell}
 \ell=\frac{g_N'}{g_N}=\frac1{50}-\frac12\log q_N<0,
 \quad z_e=B_{N,n}+eg_N A_{N,n},\quad d_e=1+eg_N,\quad
 v'=g_N(\ell A_{N,n}+A_{N,n}').
\end{equation}
The branch $u_e=|z_e|/d_e$ is active if it realizes the maximum in
\eqref{eq:U}. Both branches are retained at a tie.

\begin{lemma}[Active-branch gate]\label{lem:gate}
On the stated strip, $\sigma'_{N,+}\ge1/2$, and
\begin{equation}\label{eq:gate}
 D^+(u_e n^{-\sigma_N})\le
 \frac{n^{-\sigma_N}}{d_e}
 \left[\mathcal D(z_e,ev')-
 \left(\frac L2+\frac{eg_N'}{d_e}\right)|z_e|\right].
\end{equation}
Nonpositivity of the bracket for every active branch implies that the
corresponding summand of $V_N$ is nonincreasing. This remains valid at
numerator zeros, branch ties and the positive-part kink of $\sigma_N$.
\end{lemma}
\begin{proof}
The argument of the positive part in \eqref{eq:sigma} has negative
$y$-derivative $-3+4(1+2y)/x_N^2$. Hence its contribution to the
right slope of $\sigma_N$ is nonnegative. Differentiate $|z_e|/d_e$
using $z_e'=ev'$ and $d_e'=eg_N'$, then use $\sigma'_{N,+}\ge1/2$.
The right derivative of a maximum of two locally Lipschitz branches is
the maximum of their derivatives at the active branches. The resulting
upper-right nonpositive derivative implies monotonicity; continuity
includes the upper endpoint of the strip.
\end{proof}

\begin{proposition}[The plain gate does not imply the sharp gate]\label{prop:counter}
In the generic real signed finite-exponential-sum class, a sufficient
plain-weight Dini test does not imply the normalized sharp test.
\end{proposition}
\begin{proof}
Take $n=16$, $\sigma=\sigma_0+y/2$, $g=4^{-y}$, $B=1$, and
$A(y)=-1+3\cdot4^y-16^y$. Thus $A(0)=B$. At $y=1/2$, with $a=\log2$,
\[
 g=\tfrac12,\quad A=1,\quad A'=-4a,\quad g'=-a.
\]
The sufficient plain bracket is
\[
 g|A'|+g'|A|-\tfrac12\log n\,(|B|+g|A|)=-2a<0.
\]
Both normalized branches equal one, but their weighted right slopes,
after division by $16^{-\sigma}a$, are $2$ and $-10/3$. The maximum is
positive. The example refutes only generic automatic inheritance; it is
not an example with the theta Euler coefficients of \eqref{eq:coeff}.
\end{proof}

\section{The uniform finite-leg transfer certificate}\label{sec:transfer}

\begin{certtheorem}[Four-leg strip transfer]\label{thm:transfer}
For every leg in Table~\ref{tab:legs}, every integer $N$ in that leg, and
every $2\le n\le DN$, the function
\[
 y\longmapsto U_{g_N(y)}(B_{N,n},A_{N,n}(y))n^{-\sigma_N(y)}
\]
is nonincreasing on $[y_0,y_1]$. Consequently $V_N$ is nonincreasing.
\end{certtheorem}

\subsection{Divisor patterns and feasible integer coverage}

Put $G=\gcd(n,D)$, and order its divisors as $1=d_0<\cdots<d_{r-1}=G$.
The sector $n\le N$ has all divisors active. In sector
$d_{j-1}N<n\le d_jN$, precisely $d_j,\ldots,G$ are active. Writing
$n=Gk$ imposes $\gcd(k,D/G)=1$. If $n>GN$ there is no active divisor,
and the summand is identically zero. The numbers of nonempty sector
types are $3^{|P|}$, namely $243,81,27,9$ on the four legs.

The certificate uses coordinates $h=\log N$, $v=\log(n/N)$, $L=h+v$.
Define
\begin{equation}\label{eq:coords}
 \delta(h)=-\frac12\log\left(1-\frac{\tstar}{16}e^{-2h}\right),\quad
 \ell=\frac1{50}-h+\delta(h),\quad
 F=g_Ne^{yL}=e^{y(v+1/50+\delta(h))}.
\end{equation}
Thus the full $N$-dependence is retained. For the active divisors let
\begin{align}
 r_d(L)&=\exp\!\left[\frac{\tstar}{4}
 \left(\sum_{p\mid d}\log^2p+\log^2d-2L\log d\right)\right],\label{eq:rd}\\
 c&=\sum_{d\text{ active}}(-1)^{\omega(d)}r_d(L),
 &C_y&=\sum_{d\text{ active}}(-1)^{\omega(d)}r_d(L)d^{-y}.\label{eq:cy}
\end{align}
Then $B_{N,n}=b_{\tstar}(n)c$ and
$A_{N,n}=b_{\tstar}(n)e^{yL}C_y$. Remove the positive factor
$b_{\tstar}(n)$ and put
\begin{equation}\label{eq:QT}
 Q=FC_y,\qquad T=F[(v+1/50+\delta)C_y+\partial_yC_y].
\end{equation}
The exact branch comparison is
\begin{equation}\label{eq:activity}
 u_-\ge u_+\quad\Longleftrightarrow\quad(gc-Q)(c-gQ)\ge0.
\end{equation}
With $g=g_N$, the implemented branch predicate is nonpositivity of
\begin{equation}\label{eq:predicate}
 (1+eg)\mathcal D(c+eQ,eT)
 -\left[\frac L2(1+eg)+e\ell g\right]|c+eQ|
\end{equation}
for every possibly active branch. Multiplication by $1+eg>0$ preserves
the sign in \eqref{eq:gate}.

The interval cover embeds each actual integer pair in closed boxes for
$(h,v,y)$; sector closures only enlarge the domain. Outward
exponentiation bounds the possible integers $N,n$. The exact sector
inequalities and gcd restriction provide necessary feasibility tests.
If at most 33 possible multiples $n=Gk$ remain, the code checks all of
them against the integer leg/sector constraints and outer ratio bounds.
A cell is discarded only if these necessary conditions admit no pair.
A unique possible $n$ tightens the interval for $L$ to its logarithm.
Closed bisection children share their boundary. These operations cannot
remove a feasible pair or leave a subdivision gap.

\subsection{An equivalent predicate at absolute-value kinks}

Put
\[
 R=c-gQ,\quad S=Q-gc,\quad R'=-g(\ell Q+T),\quad S'=T-\ell gc.
\]
The identity
\[
 \max(u_-,u_+)=\frac{|R|+|S|}{1-g^2}
\]
gives the equivalent derivative numerator
\begin{equation}\label{eq:alternate}
 \mathcal D(R,R')+\mathcal D(S,S')-
 \left[\frac L2-\frac{2g^2\ell}{1-g^2}\right](|R|+|S|).
\end{equation}
A nonpositive enclosure of \eqref{eq:alternate} proves the gate even
when separate branch enclosures overestimate a tie. In either form,
a zero-containing numerator retains the upper bound on its absolute
right derivative. In \eqref{eq:predicate} the coefficient of the plus
absolute value can be negative; its complete signed interval product
is retained. No midpoint sign is used to suppress a branch.

\subsection{Analytic cases}

\begin{lemma}[A single active divisor]\label{lem:single}
Every feasible sector with exactly one active divisor satisfies the gate.
\end{lemma}
\begin{proof}
For that divisor $d$, put $m=n/d$, $a=\log m$. Then $B\ne0$,
$A=B e^{ay}$ and $0\le a\le\log N<-2\ell$. The final inequality is
$q_N/N>e^{0.04}$, which holds at the smallest $N$ and persists because
$N-\tstar/(16N)$ increases. If $a>0$ only plus is active and
\[
 \frac{d}{dy}\log u_+=\frac a2+\frac12
 \left[(\ell+a)\tanh\frac{y(\ell+a)}2
       -\ell\tanh\frac{y\ell}2\right]\le\frac a2.
\]
Indeed $x\tanh(yx/2)$ is even and increases with $|x|$, and
$|\ell+a|\le|\ell|$. Subtracting $\sigma'_{N,+}\log n$ bounds the
weighted rate by $-\tfrac12\log d\le0$. At $m=1$, both branches equal
$|B|$ identically and their weighted rates are also nonpositive.
\end{proof}

\begin{lemma}[The isolated prime zeros]\label{lem:prime}
If a full-active $G=p\in P$ cell has $n=p$ as its only possible integer,
the gate holds on every actual pair in that cell.
\end{lemma}
\begin{proof}
Here $B=0$, $A=b_{\tstar}(p)(p^y-1)>0$ and minus is uniquely active.
The weighted logarithmic rate is at most
\[
 \frac{\ell}{1-g}+\frac{\log p}{1-p^{-y}}-\frac{\log p}{2}
 \le\ell+\frac{\log p}{2}+\frac1y
 <0.52-13+1.5+\frac{1000}{187}
 =-\frac{52663}{9350}<0.
\]
Use $q_N>N^2/2$, $\log N>13$, $\log p<3$ and $y>0.187$; the
inequality $1/(1-e^{-u})\le1+1/u$ gives the middle estimate.
No other simultaneous coefficient zero is assigned this shortcut.
\end{proof}

\subsection{Execution and proof of the theorem}

The engine uses 40-digit directed Decimal arithmetic. The nearest-mode
correctly rounded exponential, logarithm and square root are padded
outward by $10^{-38}$. A cell passes by \eqref{eq:predicate},
\eqref{eq:alternate}, Lemma~\ref{lem:single} or Lemma~\ref{lem:prime};
it may also be proved empty. Otherwise it is subdivided or reported
unresolved. Table~\ref{tab:cover} gives the complete pattern partition.

\begin{table}[htbp]
\centering
\begin{tabular}{llll}
\toprule
Leg & Pattern indices & Certificate & Remaining cells\\
\midrule
0 & $0,\ldots,171$ & First pass, complete prefix & None in prefix\\
0 & $172,\ldots,242$ & Completion, full original roots & None\\
1 & $0,\ldots,80$ & First pass & None\\
2 & $0,\ldots,26$ & First pass & None\\
3 & $0,\ldots,8$ & First pass & None\\
\bottomrule
\end{tabular}
\caption{Composition of the finite-leg transfer certificate.}
\label{tab:cover}
\end{table}

\begin{proof}[Proof of Theorem~\ref{thm:transfer}]
The root boxes and necessary integer tests cover every nonempty actual
term. The first run completed legs 1--3 and the indicated prefix of
leg 0 before reaching its 30,000-box cap. Pattern 172 was only partly
evaluated then, so that work is not used for composition. The second run
restarted its full root and completed the entire suffix 172--242.
Its wrapper checks the engine hash, the first record's completed prefix
and the other complete legs, and records the first record's hash.
The exact prefix/suffix union is all 243 patterns, with no unresolved
leaf in the completed domains. Each run used at most 45 seconds and
30,000 evaluations per leg, at depth at most 36; the measured combined
interval time was 91.06 seconds and peak memory was below 22 MiB.
The records in Appendix~\ref{app:records} summarize the interval
decisions and contain accepted-leaf digests; the engine regenerates the
leaf inequalities. No conclusion is drawn from a positive
upper interval bound or a partial pattern. The sign predicates and the
analytic cases therefore prove every active gate. Apply
Lemma~\ref{lem:gate}, and add the finitely many nonincreasing summands.
\end{proof}

The first execution record alone has a false full-completion flag,
correctly reflecting its budget stop. Both records are necessary for
Theorem~\ref{thm:transfer}. A separate same-model check compares actual
integer terms with 110-digit direct arithmetic, tests that feasible toy
pairs survive pruning, and rejects the tied adverse example of
Proposition~\ref{prop:counter}. These checks examine implementation;
they do not replace the uniform interval cover or establish review
independence. No global extremal slack is inferred from the accepted
leaf thresholds or their count.

\begin{corollary}[Transfer of a positive native floor]\label{cor:floor}
If $V_N(y_0)<1$, then $\Lsharp(y)\ge\Lsharp(y_0)$ for
$y_0\le y\le y_1$.
\end{corollary}
\begin{proof}
Theorem~\ref{thm:transfer} makes $V_N$ nonincreasing. The factor $1-g_N$
is positive and increasing, and $M_N$ is nonincreasing because
$\sigma'_{N,+}\ge1/2$. Finally put $\rho_N=\tstar/[2(x_N-6)]$.
Each positive summand of $K_N$ has logarithmic rate at most
\[
 \ell+\tfrac12\log N+\rho_N\log N+1/y_0.
\]
This follows from
$\frac{d}{dy}\log(m^{\rho_Ny}-1)\le\rho_N\log m+1/y$.
For $\tstar<0.162$, $N\ge690988$,
\[
 \ell<0.52-\log N,\qquad \rho_N<0.081/N^2,
\]
so the rate is less than
$0.52-6.5+0.081/N+1000/187<-0.631$.
Thus $K_N$ decreases, and every operation in \eqref{eq:Lsharp} has
the required direction when $1-V_N(y_0)>0$.
\end{proof}

\section{A finite window and the complete approximation allowance}\label{sec:pilot}

The approximation allowance is recomputed at the new exact time; the
historical time box is not reused. A convenient statement is
\begin{proposition}[Uniform approximation budget]\label{prop:error}
For every natural cutoff $N\ge\Nzero$, $x\in W_N$ and
$y_0\le y\le y_1$ at $t=\tstar$, the inherited effective approximation
supplies
\begin{align*}
 e_A+e_B&<2.159\cdot10^{-12},&
 e_{C,0}&<2.829422\cdot10^{-7},\\
 e_A+e_B+e_{C,0}&<2.829444\cdot10^{-7}<3\cdot10^{-7}.
\end{align*}
\end{proposition}
\begin{proof}
We instantiate the audited general U1--U5 reduction \cite[Section 4]{audit4}
at $t=\tstar$, not its old fixed-parameter conclusion. Here is the
budget explicitly. Put $z_0=\tstar/(16\Nzero^2)$, $L_0=\log q_{\Nzero}$,
$x_0=x_{\Nzero}$ and
\[
 \delta_0=\frac{\tstar}{4}\frac{z_0}{1-z_0}
       +\frac{\tstar}{2x_0^2},\qquad
 \sigma_0=\frac{1+y_0}{2}+\frac{\tstar}{2}\log\Nzero-\delta_0.
\]
Use the endpoint cap of Lemma~\ref{lem:cap} below, with $m_0=2000$, to form
\begin{align*}
 S_0&=1+\sum_{n=2}^{m_0}b_{\tstar}(n)n^{-\sigma_0}
        +\Capbd_{\tstar}(m_0,\Nzero;\sigma_0),\\
 P_0&=1+e^{1/50}(1-z_0)^{-1/2}
        \exp\!\left[\frac{\tstar\log\Nzero}{2(x_0-6)}\right],\\
 U_0&=\frac{(\tstar^2/16)L_0^2+0.626}{x_0-6.66}.
\end{align*}
The error supplier bounds $e_A+e_B$ by $P_0S_0(e^{U_0}-1)$. The
$C$-term allowance is bounded by
\begin{equation}\label{eq:eCbudget}
 \exp\!\left[-\frac{1+y_0}{4}L_0-\frac{\tstar}{16}L_0^2
 +\frac{1.24(3^{y_1}+3^{-y_1})}{\Nzero-1/8}
 +\frac{3\sqrt{L_0^2+\pi^2/4}+10.50}{x_0-12}\right].
\end{equation}
The head/cap monotonicity and domain gates of U1--U5 are rechecked in
\path{TAIL_T157_M600000_RECORD.json}; its two U3 sign gates are
less than $-0.42558$ and $-0.47805$. Directed evaluation gives the
respective upper bounds $2.158316897\cdot10^{-12}$ and
$2.829421823\cdot10^{-7}$. The reduction is uniform for all $N\ge\Nzero$
and the full new strip. The error formulas are imports from
\cite{polymath,audit4}; the parameter substitution and enclosure are
recorded in \cite{note01} and Appendix~\ref{app:records}.
\end{proof}

\begin{certtheorem}[First-window strip margin]\label{thm:pilot}
At $t=\tstar$, throughout $x\in W_{690988}$ and $y_0\le y\le y_1$,
\begin{equation}\label{eq:pilotmargin}
 \left|\frac{H_{\tstar}(x+iy)}{B_{\tstar}(x+iy)}\right|
 >0.0017752269101384707.
\end{equation}
\end{certtheorem}
\begin{proof}
The outward fixed-point pilot encloses
\begin{align*}
 0.0017755269101384707588
 &< L_{690988}^{\mathrm{sh}}(y_0)\\
 &<0.0017755269101384708095.
\end{align*}
It uses the exact enumeration $n=Gk$, $\gcd(k,D/G)=1$, signed
convolutions and the weight $W_g$. In the record, $S$ denotes the
\emph{unnormalized} mass $(1-g)V_N$, and $Q=1-g-S$; this $S$ must not
be mistaken for $V_N$. The 110-bit fixed-point sum allowance is below
$4.54\cdot10^{-20}$, and the scalar/table supplier uses inherited
45-digit directed arithmetic. Error propagation uses the 1-Lipschitz
absolute-value and maximum maps, including the multiplier
$(1-g)/(1+g)$ and directly enclosed $G^{-\sigma_N}$ scalars.
The record's positive numerator implies $V_N(y_0)<1$.
Corollary~\ref{cor:floor} and Proposition~\ref{prop:binding} transfer this
floor through the entire strip and window. Subtract the full
$3\cdot10^{-7}$ allowance of Proposition~\ref{prop:error} using
\eqref{eq:approx}.
\end{proof}

\section{An infinite-tail estimate with a larger head}\label{sec:tail}

This section re-instantiates the \emph{generic arguments} of the preceding
tail audit \cite[Sections 3.1--3.8]{audit4}; its old stated time interval
does not contain $\tstar$. The partition and reduction are given here,
and every new application gate is reevaluated.

\begin{lemma}[Endpoint cap]\label{lem:cap}
For $t\ge0$, integers $1\le a<N$ and $\sigma>(t/2)\log N$,
\[
 \sum_{a<n\le N}b_t(n)n^{-\sigma}
 \le\Capbd_t(a,N;\sigma)
 :=\log(N/a)\max_{u\in\{a,N\}}e^{(1-\sigma)\log u+(t/4)\log^2u}.
\]
\end{lemma}
\begin{proof}
The summand decreases, so its sum is bounded by the integral from $a$
to $N$. With $u=e^v$, the integrand is
$\exp((1-\sigma)v+tv^2/4)$, a convex exponent whose maximum on the
logarithmic interval is at an endpoint. This is the subrange version of
Polymath's Lemma 8.2 \cite{polymath}.
\end{proof}

The tail uses a different rational mollifier. Put
$\mathcal S=\{1,2,3,4,5,6,7,10,11,13,14\}$ and, in that order, set
\begin{equation}\label{eq:tailnu}
 (\nu_d)_{d\in\mathcal S}=
 \left(1,-\frac{1021}{1000},-\frac{1054}{1000},-\frac9{200},
 -\frac{1119}{1000},\frac{1001}{1000},-\frac{1043}{1000},
 \frac{128}{125},-\frac{161}{200},-\frac{447}{500},
 \frac{456373}{500000}\right).
\end{equation}
These coefficients are not the finite Euler coefficients $\lambda_d$.
Let $\mathcal E_\nu(s)=\sum_{d\in\mathcal S}\nu_dd^{-s}$ and fix
$M=600000<\Nstar$. With $c_m=\sum_{d\mid m,\,d\in\mathcal S}
\nu_db_{\tstar}(m/d)$, the exact pair partition is
\begin{align}\label{eq:partition}
 \mathcal E_\nu(s)\sum_{k\le N}b_{\tstar}(k)k^{-s}
 &=1+\sum_{m=2}^M c_mm^{-s}\\
 &\quad+\sum_{d\in\mathcal S}\nu_dd^{-s}
     \sum_{\lfloor M/d\rfloor<k\le N}b_{\tstar}(k)k^{-s}.\notag
\end{align}
All pairs with $dk\le M$ occur in the head, since $M<N$; the rest
occur in exactly one routed sum. The historical nonnegative term
\texttt{OV} bounds part of this routed remainder a second time and is
not needed in this partition.

Let $x_*=x_{\Nstar}$ in this section only, and put
\begin{align}
 z_*&=\frac{\tstar}{16\Nstar^2},&
 \widehat\delta&=\frac{\tstar}{4}\frac{z_*}{1-z_*}
                          +\frac{\tstar}{2x_*^2},&
 \widehat k&=\frac{\tstar}{2(x_*-6)},\label{eq:tailfreeze}\\
 \sigma_1(N,y)&=\frac{1+y}{2}+\frac{\tstar}{2}\log N-\widehat\delta,
 &\sigma_2(N,y)&=\sigma_1(N,y)-y-\widehat k,\label{eq:sig12}\\
 G(N,y)&=\exp\!\left(\frac y{50}-\frac y2\log(N^2-\tstar/16)\right).
 &&\label{eq:tailG}
\end{align}
The imported approximation bounds imply $\Rea s_*\ge\sigma_1$,
$|\kappa|\le\widehat k$ and $|\gamma|\le G$ for every $N\ge\Nstar$.
Here $-\log(1-z)\le z/(1-z)$ and the positive part in
\eqref{eq:sigma} is bounded by one. Define
\begin{align}
 P&=\sum_{m=2}^M|c_m|m^{-\sigma_1},\qquad
 \mathfrak M=\sum_{d\in\mathcal S}|\nu_d|d^{-\sigma_1},\label{eq:Ptail}\\
 T_R&=\sum_{d\in\mathcal S}|\nu_d|d^{-\sigma_1}
       \Capbd_{\tstar}(\lfloor M/d\rfloor,N;\sigma_1),\label{eq:TR}\\
 A_B&=G\left[\sum_{n=1}^M b_{\tstar}(n)n^{-\sigma_2}
                  +\Capbd_{\tstar}(M,N;\sigma_2)\right],\label{eq:AB}\\
 D'&=P+T_R+\mathfrak M A_B.\label{eq:Dprime}
\end{align}
The partition and ordinary triangle inequality give
\begin{equation}\label{eq:tailbinding}
 |\mathcal E_\nu(s_*)f_{\tstar}-1|\le D',\qquad
 |f_{\tstar}|\ge\frac{1-D'}{\mathfrak M}\quad(D'<1).
\end{equation}
The second inequality follows because $|\mathcal E_\nu|\le\mathfrak M$;
$D'<1$ itself ensures that the mollifier and the product are nonzero.

\begin{proposition}[Order of the tail reductions]\label{prop:tailorder}
The majorant $D'(N,y)$ is nonincreasing first in $y$, then in $N$ at the
lower $y$ endpoint, provided the recorded cap gates below hold.
The contraction's $A$-cap all-index gate is needed only near $y_0$;
the error caps still require the full strip hull.
\end{proposition}
\begin{proof}
For fixed $N$, all $B$-side terms decrease with $y$ through
$\partial_y\sigma_1=1/2$. Every $G$-weighted $A$-side head term or cap
endpoint $u\le N$ has logarithmic $y$-rate at most
\[
 Y_M(N)=\frac1{50}-\frac12\log(N^2-\tstar/16)+\frac12\log N<0.
\]
The logarithmic cap width is independent of $y$. Thus $D'$ decreases
throughout the full strip, before any $N$ reduction.

For a fixed cap left endpoint $a$, set $q=\log N$, $r=\log a$, and
$\sigma(N)=\beta+(\tstar/2)q$. The two cap candidates have logarithmic
$q$-derivatives
\[
 \frac1{q-r}-\frac{\tstar}{2}r,\qquad
 \frac1{q-r}+1-\sigma(N).
\]
Both are negative if
\begin{equation}\label{eq:Ngates}
 \frac{\tstar}{2}r(q-r)>1,\qquad
 (\sigma(N)-1)(q-r)>1.
\end{equation}
The left sides increase with $q$ once they hold. At the reduced lower
$y$ value these gates therefore make all routed and $A$ caps decrease
with $N$. The heads and mollifier majorant decrease as their exponents
increase; $G$ also decreases. This proves the contraction reduction.
The approximation-error supplier does not use this $y$ reduction, so
its two error caps must retain the full-hull all-index gates.
\end{proof}

\begin{certtheorem}[Tail at the unchanged starting index]\label{thm:tail}
For $t=\tstar$, every $N\ge\Nstar$, $x\in W_N$ and $y_0\le y\le y_1$,
\begin{equation}\label{eq:tailmargin}
 \left|H_{\tstar}(x+iy)/B_{\tstar}(x+iy)\right|>0.02303064.
\end{equation}
\end{certtheorem}
\begin{proof}
The contraction is enclosed at $N=\Nstar$ on
\[
 I_{\rm box}=[1872719/10^7,23409/125000],
\]
containing $y_0$. The interval's lower endpoint lies below $y_0$, so
Proposition~\ref{prop:tailorder} reduces the required strip to that box.
The error supplier retains
\[
 I_{\rm ext}=[1872719/10^7,8282512/10^7],
\]
which covers $y_1$; the historical top $0.8231039$ would not suffice.
Directed evaluation at the exact new time gives Table~\ref{tab:tail}.
The least routed cap gates exceed $1.93875$ and $1.45491$; the
$A$-cap second gate on $I_{\rm box}$ exceeds $1.10728$. Its historical
extended-hull test does not pass (its computed lower endpoint is about
$0.51236$); that test is not needed after the stated reduction order.
Both full-hull error caps pass, with least gate exceeding $1.97475$.
The validity gates $\sigma_j>(\tstar/2)\log N$ and all remaining domain
and monotonicity gates are retained in the record.

The exact head is $M=600000$; the two error heads use 3000 terms. Their
inherited error formulas, including the $10.50$ weld, give the complete
tail allowance in Table~\ref{tab:tail}. Equations
\eqref{eq:approx} and \eqref{eq:tailbinding}, and the all-index reductions,
then prove the result. All numerical comparisons are reevaluations of
the generic audited arguments at $\tstar$, not applications of the old
fixed-time theorem outside its domain.
\end{proof}

\begin{table}[htbp]
\centering
\begin{tabular}{ll}
\toprule
Quantity at the declared hull & Safe outward bound\\
\midrule
$P$ & $<0.214842556622$\\
$T_R$ & $<0.214622245029$\\
$\mathfrak M$ & $<1.631076048739$\\
$A_B$ & $<0.326760019216$\\
$D'$ & $<0.962435242680$\\
$e_A+e_B+e_{C,0}$ in the tail & $<1.4900275280\cdot10^{-8}$\\
$(1-D')/\mathfrak M-(e_A+e_B+e_{C,0})$
 & $>0.02303064473$\\
\bottomrule
\end{tabular}
\caption{New-time tail budget with $M=600000$. The final row uses the
full-precision enclosures, rather than recombination of rounded table entries.}
\label{tab:tail}
\end{table}

At the unchanged head $153814$, the same new-time budget has an upper
bound about $1.15923637$ for $D'$ and does not close. This is a failure of
that head and majorant to certify contraction, not a theorem of
nonvanishing failure or an obstruction to other heads. Similarly, retaining
the redundant \texttt{OV} padding in the successful case would prevent
this particular budget from closing. The raw historical checks and their
replacement reasons are preserved; adverse required gates are not removed.

\section{The remaining finite bridge and the conditional target}\label{sec:remaining}

Polymath's Theorem 1.2 \cite{polymath} converts three inputs into
$\Lambda\le t_0+y_0^2/2$: initial verification to height $X/2$,
final-time nonvanishing for $x\ge X+\sqrt{1-y_0^2}$ and
$y_0\le y\le\sqrt{1-2t_0}$, and an intermediate curved barrier.
For
\[
 X=6000000185827,
\]
the Platt--Trudgian verified height \cite{platt} covers $X/2$.
The new curved barrier at $t_0=\tstar$ is contained in the historical
closed box
\[
 [X,X+1]\times[0.1809,1]\times[0,0.16125].
\]
Indeed its horizontal width is less than one, its lower ordinate is at
least $y_0>0.1809$, and $\tstar<0.16125$. The historical barrier is
inherited on the status recorded by its audit \cite{audit5}; it is not
regenerated here. The first natural window still contains
$X+\sqrt{1-y_0^2}$, and the half-open $W_N$ tile without gaps.

\begin{corollary}[Explicit remaining obligation]\label{cor:conditional}
Assume the imported initial-verification and historical barrier inputs
just described. If a valid new floor certificate establishes
\[
 L_N^{\mathrm{sh}}(y_0)>3\cdot10^{-7}
 \quad\text{for every integer }690989\le N<3840000
\]
with its leg's coefficients and all producer or amortization errors paid,
then $\Lambda\le0.1745354$.
\end{corollary}
\begin{proof}
The first floor is Theorem~\ref{thm:pilot}; the remaining finite floors
transfer by Corollary~\ref{cor:floor} and pay
Proposition~\ref{prop:error}. Theorem~\ref{thm:tail} supplies the infinite
exterior. These cover the entire final-time region. Apply the imported
criterion and compute
$\tstar+y_0^2/2=872677/5000000=0.1745354$.
\end{proof}

The premise about the remaining floors is open. Neither the historical
plain-weight rows at time $0.16125$, the favorable sampled sharp rows,
nor the transfer theorem proves it. No monotonicity in $N$ is assumed.
This draft therefore reports a completed transfer calculation, one
finite-window margin and a revised infinite exterior, not a complete
certificate of the conditional conclusion.

The control-family screen is also quantitative. The generic paired
algebra has no shape premise that excludes colliding sources. Euler-factor
controls change the coefficients and conductor; weighted-square,
reflected, positive-primitive and earlier positive-source controls do not
provide the complete theta approximation with its stated error allowance
\cite{note01}. Their own representations and errors would have to replace
these inputs. A valid method may certify a collision-free box of a
control, but cannot produce a positive complete margin on a region
containing its collision. The present theorem makes no such claim.

\appendix
\section{Evidence and reproduction}\label{app:records}

All current evidence below is under
\path{papers/investigations-3/02_sharp_finite_sum_bounds/numerics/round-1/}
in the research checkout. The source notes are cited as
\cite{note01,note02}. Files and imported helper pins are bound by the
build record accompanying this manuscript.

\begin{table}[htbp]
\small
\begin{tabularx}{\textwidth}{@{}p{0.40\textwidth}X@{}}
\toprule
File & Role\\
\midrule
\path{certify_sharp_pilot.py} & Outward single-row implementation; one native floor.\\
\path{SHARP_PILOT_RECORD_20261017.json} & First floor and fixed-point/scalar allowances.\\
\path{check_sharp_y_transfer_gate.py} & Actual-integer domain cover and both derivative predicates.\\
\path{SHARP_Y_TRANSFER_GATE_RECORD_20261017.json} & Three complete legs and the largest-leg prefix.\\
\path{complete_sharp_y_transfer_gate.py} & Checks the complete prefix and starts the original suffix roots.\\
\path{SHARP_Y_TRANSFER_COMPLETION_RECORD_20261017.json} & Completed suffix and exact union; required with the preceding record.\\
\path{check_gate_arithmetic.py} & Separate same-model implementation checks; not uniform coverage.\\
\path{check_t157_tail_budget.py} & Hash-pinned generic error/tail adapter with explicit substitutions.\\
\path{TAIL_T157_M600000_RECORD.json} & Uniform error allowance and successful new-time tail.\\
\path{TAIL_T157_M153814_RECORD.json} & Inherited-head budget diagnostic.\\
\bottomrule
\end{tabularx}
\caption{The small evidence package. Floating scout values are diagnostics only.}
\label{tab:evidence}
\end{table}

The transfer engine checks its interval primitives, exact rational
one-sided derivative identities, and toy integer membership. The
additional checker compares 1080 actual term points at 110 digits,
tests 171 feasible toy pairs through 2223 containing child cells, and
rejects the adverse tied control. These counts account for checks; the
proof uses analytic arguments and the interval signs, not the number of
checks. The tail supplier uses 50-digit directed arithmetic with
$10^{-48}$ transcendental padding, and explicitly retains failed legacy
comparisons whose historical target, floor or reduction has changed.

From the numerical directory
\path{papers/investigations-3/02_sharp_finite_sum_bounds/numerics/round-1/},
a cheap implementation check is
\begin{quote}\small\ttfamily
python3 check\_gate\_arithmetic.py --output /tmp/sharp-gate-check.json
\end{quote}
The exact bounded-run commands and plans are in the numerical README.
The completion wrapper uses the preserved first record beside the
engine; it does not automatically consume an arbitrary freshly generated
first pass. Resource stops remain incomplete and must be reported.
Each source file and record is below 1 MiB; no large row corpus is
created. Hashes establish input and execution identity, while arithmetic
and the stated analytic reductions establish the inequalities.

\section*{Acknowledgements and assistance}
This manuscript was drafted for Edward Baker with substantial LLM
assistance from OpenAI's \texttt{gpt-6.1-sol} in Codex, reasoning effort
\texttt{ultra}, including same-model subagents. Assistance covered
mathematical derivation, code authoring, numerical implementation checks
and exposition. Compilation and layout inspection are authoring checks;
they do not constitute the deferred round-level or specialist review.

\begin{thebibliography}{9}
\bibitem{polymath}
D. H. J. Polymath,
\emph{Effective approximation of heat flow evolution of the Riemann
$\xi$ function, and a new upper bound for the de Bruijn--Newman constant},
arXiv:1904.12438v2 (2019).
\url{https://arxiv.org/abs/1904.12438v2}.

\bibitem{platt}
D. J. Platt and T. S. Trudgian,
\emph{The Riemann hypothesis is true up to $3\cdot10^{12}$},
arXiv:2004.09765 (2020).
\url{https://arxiv.org/abs/2004.09765}.

\bibitem{note01}
Prepared for E. Baker with LLM assistance,
\emph{Sharp finite sum contract and first feasibility budget},
17 October 2026, \texttt{gpt-6.1-sol}, effort \texttt{ultra}.
Investigation 02, round-one Note 01,
\path{../../notes/round-1/01_SHARP_WEIGHT_CONTRACT_AND_FEASIBILITY_20261017.md}.

\bibitem{note02}
Prepared for E. Baker with LLM assistance,
\emph{Bounded implementation of the sharp strip transfer gate},
17 October 2026, \texttt{gpt-6.1-sol}, effort \texttt{ultra}.
Investigation 02, round-one Note 02,
\path{../../notes/round-1/02_BOUNDED_ACTIVE_BRANCH_TRANSFER_IMPLEMENTATION_20261017.md}.

\bibitem{sharpreview}
Internal LLM audit,
\emph{Cross-model review of the first Newman certificate audit session},
12 October 2026, Section 3.3.
\path{papers/investigations-2/newman_certificate_audit/reviews/1_CROSS_MODEL_REVIEW_OF_NOTES_1_2_20261012.md}.

\bibitem{audit3}
Internal LLM audit,
\emph{Certified row replay and audit of the finite-lane reductions},
12 October 2026.
\path{papers/investigations-2/newman_certificate_audit/notes/3_CERTIFIED_ROW_REPLAY_AND_FINITE_LANE_REDUCTIONS_20261012.md}.

\bibitem{audit4}
Internal LLM audit,
\emph{Tail lemma and uniform error budget: audit of the remaining analytic
suppliers of hypothesis (ii)}, 12 October 2026.
\path{papers/investigations-2/newman_certificate_audit/notes/4_TAIL_LEMMA_AND_UNIFORM_ERROR_BUDGET_AUDIT_20261012.md}.

\bibitem{audit5}
Internal LLM audit,
\emph{Barrier certificate audit}, 13 October 2026.
\path{papers/investigations-2/newman_certificate_audit/notes/5_BARRIER_CERTIFICATE_AUDIT_20261013.md}.
\end{thebibliography}
\end{document}
